\section{Rings of finite global dimension}
\label{section-ring-finite-gl-dim}

\noindent
The following lemma is often used to compare different
projective resolutions of a given module.

\begin{lemma}[Schanuel's lemma]
\label{lemma-Schanuel}
Let $R$ be a ring. Let $M$ be an $R$-module.
Suppose that
$$
0 \to K \xrightarrow{c_1} P_1 \xrightarrow{p_1} M \to 0
\quad\text{and}\quad
0 \to L \xrightarrow{c_2} P_2 \xrightarrow{p_2} M \to 0
$$
are two short exact sequences, with $P_i$ projective.
Then $K \oplus P_2 \cong L \oplus P_1$. More precisely,
there exist a commutative diagram
$$
\xymatrix{
0 \ar[r] &
K \oplus P_2 \ar[r]_{(c_1, \text{id})} \ar[d] &
P_1 \oplus P_2 \ar[r]_{(p_1, 0)} \ar[d] &
M \ar[r] \ar@{=}[d] &
0 \\
0 \ar[r] &
P_1 \oplus L \ar[r]^{(\text{id}, c_2)} &
P_1 \oplus P_2 \ar[r]^{(0, p_2)} &
M \ar[r] &
0
}
$$
whose vertical arrows are isomorphisms.
\end{lemma}

\begin{proof}
Consider the module $N$ defined by the short exact sequence
$0 \to N \to P_1 \oplus P_2 \to M \to 0$,
where the last map is the sum of the two maps
$P_i \to M$. It is easy to see that the projection
$N \to P_1$ is surjective with kernel $L$, and that
$N \to P_2$ is surjective with kernel $K$.
Since $P_i$ are projective we have $N \cong K \oplus P_2
\cong L \oplus P_1$. This proves the first statement.

\medskip\noindent
To prove the second statement (and to reprove the first), choose
$a : P_1 \to P_2$ and $b : P_2 \to P_1$ such that
$p_1 = p_2 \circ a$ and $p_2 = p_1 \circ b$. This is possible
because $P_1$ and $P_2$ are projective. Then we get a commutative diagram
$$
\xymatrix{
0 \ar[r] &
K \oplus P_2 \ar[r]_{(c_1, \text{id})} &
P_1 \oplus P_2 \ar[r]_{(p_1, 0)} &
M \ar[r] &
0 \\
0 \ar[r] &
N \ar[r] \ar[d] \ar[u] &
P_1 \oplus P_2 \ar[r]_{(p_1, p_2)}
\ar[d]_S \ar[u]^T &
M \ar[r] \ar@{=}[d] \ar@{=}[u] &
0 \\
0 \ar[r] &
P_1 \oplus L \ar[r]^{(\text{id}, c_2)} &
P_1 \oplus P_2 \ar[r]^{(0, p_2)} &
M \ar[r] &
0
}
$$
with $T$ and $S$ given by the matrices
$$
S = \left(
\begin{matrix}
\text{id} & 0 \\
a & \text{id}
\end{matrix}
\right)
\quad\text{and}\quad
T = \left(
\begin{matrix}
\text{id} & b \\
0 & \text{id}
\end{matrix}
\right)
$$
Then $S$, $T$ and the maps $N \to P_1 \oplus L$
and $N \to K \oplus P_2$ are isomorphisms as desired.
\end{proof}

\begin{definition}
\label{definition-finite-proj-dim}
Let $R$ be a ring. Let $M$ be an $R$-module. We say $M$ has
{\it finite projective dimension} if it has a finite length
resolution by projective $R$-modules. The minimal length of such a
resolution is called the {\it projective dimension}
of $M$.
\end{definition}

\noindent
It is clear that the projective dimension of $M$ is $0$ if and
only if $M$ is a projective module.
The following lemma explains to what extent the projective
dimension is independent of the choice of a projective
resolution.

\begin{lemma}
\label{lemma-independent-resolution}
Let $R$ be a ring. Suppose that $M$ is an $R$-module of projective
dimension $d$. Suppose that $F_e \to F_{e-1} \to \ldots \to F_0 \to M \to 0$
is exact with $F_i$ projective and $e \geq d - 1$.
Then the kernel of $F_e \to F_{e-1}$ is projective
(or the kernel of $F_0 \to M$ is projective in case
$e = 0$).
\end{lemma}

\begin{proof}
We prove this by induction on $d$. If $d = 0$, then
$M$ is projective. In this case there is a splitting
$F_0 = \Ker(F_0 \to M) \oplus M$, and hence
$\Ker(F_0 \to M)$ is projective. This finishes
the proof if $e = 0$, and if $e > 0$, then replacing
$M$ by $\Ker(F_0 \to M)$ we decrease $e$.

\medskip\noindent
Next assume $d > 0$.
Let $0 \to P_d \to P_{d-1} \to \ldots \to P_0 \to M \to 0$
be a minimal length finite resolution with $P_i$ projective.
According to
Schanuel's Lemma \ref{lemma-Schanuel}
we have
$P_0 \oplus \Ker(F_0 \to M) \cong F_0 \oplus \Ker(P_0 \to M)$.
This proves the case $d = 1$, $e = 0$, because then the right
hand side is $F_0 \oplus P_1$ which is projective. Hence now we may
assume $e > 0$. The module $F_0 \oplus \Ker(P_0 \to M)$ has the
finite projective resolution
$$
0 \to P_d \to P_{d-1} \to \ldots \to
P_2 \to P_1 \oplus F_0 \to \Ker(P_0 \to M) \oplus F_0 \to 0
$$
of length $d - 1$. By induction applied to the exact sequence
$$
F_e \to F_{e-1} \to \ldots \to F_2 \to P_0 \oplus F_1 \to
P_0 \oplus \Ker(F_0 \to M) \to 0
$$
of length $e - 1$ we conclude $\Ker(F_e \to F_{e - 1})$
is projective (if $e \geq 2$)
or that $\Ker(F_1 \oplus P_0 \to F_0 \oplus P_0)$ is projective.
This implies the lemma.
\end{proof}

\begin{lemma}
\label{lemma-what-kind-of-resolutions}
Let $R$ be a ring. Let $M$ be an $R$-module. Let $d \geq 0$.
The following are equivalent
\begin{enumerate}
\item $M$ has projective dimension $\leq d$,
\item there exists a resolution
$0 \to P_d \to P_{d - 1} \to \ldots \to P_0 \to M \to 0$
with $P_i$ projective,
\item for some resolution
$\ldots \to P_2 \to P_1 \to P_0 \to M \to 0$ with
$P_i$ projective we have $\Ker(P_{d - 1} \to P_{d - 2})$
is projective if $d \geq 2$, or $\Ker(P_0 \to M)$ is projective if
$d = 1$, or $M$ is projective if $d = 0$,
\item for any resolution
$\ldots \to P_2 \to P_1 \to P_0 \to M \to 0$ with
$P_i$ projective we have $\Ker(P_{d - 1} \to P_{d - 2})$
is projective if $d \geq 2$, or $\Ker(P_0 \to M)$ is projective if
$d = 1$, or $M$ is projective if $d = 0$.
\end{enumerate}
\end{lemma}

\begin{proof}
The equivalence of (1) and (2) is the definition of projective
dimension, see Definition \ref{definition-finite-proj-dim}.
We have (2) $\Rightarrow$ (4) by Lemma \ref{lemma-independent-resolution}.
The implications (4) $\Rightarrow$ (3) and (3) $\Rightarrow$ (2) are
immediate.
\end{proof}

\begin{lemma}
\label{lemma-what-kind-of-resolutions-local}
Let $R$ be a local ring. Let $M$ be an $R$-module. Let $d \geq 0$.
The equivalent conditions (1) -- (4) of
Lemma \ref{lemma-what-kind-of-resolutions}
are also equivalent to
\begin{enumerate}
\item[(5)] there exists a resolution
$0 \to P_d \to P_{d - 1} \to \ldots \to P_0 \to M \to 0$
with $P_i$ free.
\end{enumerate}
\end{lemma}

\begin{proof}
Follows from Lemma \ref{lemma-what-kind-of-resolutions} and
Theorem \ref{theorem-projective-free-over-local-ring}.
\end{proof}

\begin{lemma}
\label{lemma-what-kind-of-resolutions-Noetherian}
Let $R$ be a Noetherian ring. Let $M$ be a finite $R$-module.
Let $d \geq 0$. The equivalent conditions (1) -- (4) of
Lemma \ref{lemma-what-kind-of-resolutions}
are also equivalent to
\begin{enumerate}
\item[(6)] there exists a resolution
$0 \to P_d \to P_{d - 1} \to \ldots \to P_0 \to M \to 0$
with $P_i$ finite projective.
\end{enumerate}
\end{lemma}

\begin{proof}
Choose a resolution $\ldots \to F_2 \to F_1 \to F_0 \to M \to 0$
with $F_i$ finite free (Lemma \ref{lemma-resolution-by-finite-free}).
By Lemma \ref{lemma-what-kind-of-resolutions} we see that
$P_d = \Ker(F_{d - 1} \to F_{d - 2})$ is projective at least if $d \geq 2$.
Then $P_d$ is a finite $R$-module as $R$ is Noetherian and
$P_d \subset F_{d - 1}$ which is finite free.
Whence $0 \to P_d \to F_{d - 1} \to \ldots \to F_1 \to F_0 \to M \to 0$
is the desired resolution.
\end{proof}

\begin{lemma}
\label{lemma-what-kind-of-resolutions-Noetherian-local}
Let $R$ be a local Noetherian ring. Let $M$ be a finite $R$-module.
Let $d \geq 0$. The equivalent conditions (1) -- (4) of
Lemma \ref{lemma-what-kind-of-resolutions},
condition (5) of Lemma \ref{lemma-what-kind-of-resolutions-local},
and condition (6) of Lemma \ref{lemma-what-kind-of-resolutions-Noetherian}
are also equivalent to
\begin{enumerate}
\item[(7)] there exists a resolution
$0 \to F_d \to F_{d - 1} \to \ldots \to F_0 \to M \to 0$
with $F_i$ finite free.
\end{enumerate}
\end{lemma}

\begin{proof}
This follows from Lemmas \ref{lemma-what-kind-of-resolutions},
\ref{lemma-what-kind-of-resolutions-local}, and
\ref{lemma-what-kind-of-resolutions-Noetherian}
and because a finite projective module over a local ring
is finite free, see Lemma \ref{lemma-finite-projective}.
\end{proof}

\begin{lemma}
\label{lemma-projective-dimension-ext}
Let $R$ be a ring. Let $M$ be an $R$-module. Let $n \geq 0$.
The following are equivalent
\begin{enumerate}
\item $M$ has projective dimension $\leq n$,
\item $\Ext^i_R(M, N) = 0$ for all $R$-modules $N$ and all
$i \geq n + 1$, and
\item $\Ext^{n + 1}_R(M, N) = 0$ for all $R$-modules $N$.
\end{enumerate}
\end{lemma}

\begin{proof}
Assume (1). Choose a free resolution $F_\bullet \to M$ of $M$. Denote
$d_e : F_e \to F_{e - 1}$. By
Lemma \ref{lemma-independent-resolution}
we see that $P_e = \Ker(d_e)$ is projective for $e \geq n - 1$.
This implies that $F_e \cong P_e \oplus P_{e - 1}$ for $e \geq n$
where $d_e$ maps the summand $P_{e - 1}$ isomorphically to $P_{e - 1}$
in $F_{e - 1}$. Hence, for any $R$-module $N$ the complex
$\Hom_R(F_\bullet, N)$ is split exact in degrees $\geq n + 1$.
Whence (2) holds. The implication (2) $\Rightarrow$ (3) is trivial.

\medskip\noindent
Assume (3) holds. If $n = 0$ then $M$ is projective by
Lemma \ref{lemma-characterize-projective}
and we see that (1) holds. If $n > 0$ choose a free $R$-module $F$
and a surjection $F \to M$ with kernel $K$. By
Lemma \ref{lemma-reverse-long-exact-seq-ext}
and the vanishing of $\Ext_R^i(F, N)$ for all $i > 0$ by part (1)
we see that $\Ext_R^n(K, N) = 0$ for all $R$-modules $N$.
Hence by induction we see that $K$ has projective dimension $\leq n - 1$.
Then $M$ has projective dimension $\leq n$ as any finite projective
resolution of $K$ gives a projective resolution of length one more
for $M$ by adding $F$ to the front.
\end{proof}

\begin{lemma}
\label{lemma-exact-sequence-projective-dimension}
Let $R$ be a ring. Let $0 \to M' \to M \to M'' \to 0$ be a short
exact sequence of $R$-modules.
\begin{enumerate}
\item If $M$ has projective dimension $\leq n$ and $M''$
has projective dimension $\leq n + 1$, then $M'$ has projective
dimension $\leq n$.
\item If $M'$ and $M''$ have projective dimension
$\leq n$ then $M$ has projective dimension $\leq n$.
\item If $M'$ has projective dimension $\leq n$ and
$M$ has projective dimension $\leq n + 1$ then
$M''$ has projective dimension $\leq n + 1$.
\end{enumerate}
\end{lemma}

\begin{proof}
Combine the characterization of projective dimension in
Lemma \ref{lemma-projective-dimension-ext}
with the long exact sequence of ext groups in
Lemma \ref{lemma-reverse-long-exact-seq-ext}.
\end{proof}

\begin{definition}
\label{definition-finite-gl-dim}
Let $R$ be a ring. The ring
$R$ is said to have {\it finite global dimension}
if there exists an integer $n$ such that
every $R$-module has a resolution by
projective $R$-modules of length at most $n$.
The minimal such $n$ is then called the {\it global dimension}
of $R$.
\end{definition}

\noindent
The argument in the proof of the following lemma can be found
in the paper \cite{Auslander} by Auslander.

\begin{lemma}
\label{lemma-colimit-projective-dimension}
Let $R$ be a ring. Suppose we have a module $M = \bigcup_{e \in E} M_e$
where the $M_e$ are submodules well-ordered by inclusion. Assume the quotients
$M_e/\bigcup\nolimits_{e' < e} M_{e'}$ have projective dimension $\leq n$.
Then $M$ has projective dimension $\leq n$.
\end{lemma}

\begin{proof}
We will prove this by induction on $n$.

\medskip\noindent
Base case: $n = 0$. Then $P_e = M_e/\bigcup_{e' < e} M_{e'}$ is projective.
Thus we may choose a section $P_e \to M_e$ of the projection $M_e \to P_e$.
We claim that the induced map $\psi : \bigoplus_{e \in E} P_e \to M$ is an
isomorphism. Namely, if $x = \sum x_e \in \bigoplus P_e$ is nonzero,
then we let $e_{max}$ be maximal such that $x_{e_{max}}$ is nonzero
and we conclude that $y = \psi(x) = \psi(\sum x_e)$ is nonzero because
$y \in M_{e_{max}}$ has nonzero image $x_{e_{max}}$ in $P_{e_{max}}$.
On the other hand, let $y \in M$. Then $y \in M_e$ for some $e$.
We show that $y \in \Im(\psi)$ by transfinite induction on $e$.
Let $x_e \in P_e$ be the image of $y$. Then
$y - \psi(x_e) \in \bigcup_{e' < e} M_{e'}$.
By induction hypothesis we conclude that $y - \psi(x_e) \in \Im(\psi)$
hence $y \in \Im(\psi)$. Thus the claim is true and
$\psi$ is an isomorphism. We conclude that $M$ is projective as
a direct sum of projectives, see
Lemma \ref{lemma-direct-sum-projective}.

\medskip\noindent
If $n > 0$, then for $e \in E$ we denote $F_e$ the free $R$-module
on the set of elements of $M_e$. Then we have a system of
short exact sequences
$$
0 \to K_e \to F_e \to M_e \to 0
$$
over the well-ordered set $E$. Note that the transition maps
$F_{e'} \to F_e$ and $K_{e'} \to K_e$ are injective too.
Set $F = \bigcup F_e$ and $K = \bigcup K_e$. Then
$$
0 \to
K_e/\bigcup\nolimits_{e' < e} K_{e'} \to
F_e/\bigcup\nolimits_{e' < e} F_{e'} \to
M_e/\bigcup\nolimits_{e' < e} M_{e'} \to 0
$$
is a short exact sequence of $R$-modules too and
$F_e/\bigcup_{e' < e} F_{e'}$ is the free $R$-module on the
set of elements in $M_e$ which are not contained in $\bigcup_{e' < e} M_{e'}$.
Hence by
Lemma \ref{lemma-exact-sequence-projective-dimension}
we see that the projective dimension of $K_e/\bigcup_{e' < e} K_{e'}$
is at most $n - 1$. By induction we conclude that $K$ has projective
dimension at most $n - 1$. Whence $M$ has projective dimension at most
$n$ and we win.
\end{proof}

\begin{lemma}
\label{lemma-finite-gl-dim}
Let $R$ be a ring. The following are equivalent
\begin{enumerate}
\item $R$ has finite global dimension $\leq n$,
\item every finite $R$-module has projective dimension $\leq n$, and
\item every cyclic $R$-module $R/I$ has projective dimension $\leq n$.
\end{enumerate}
\end{lemma}

\begin{proof}
It is clear that (1) $\Rightarrow$ (2) and (2) $\Rightarrow$ (3).
Assume (3). Choose a set $E \subset M$ of generators of $M$.
Choose a well ordering on $E$. For $e \in E$ denote
$M_e$ the submodule of $M$ generated by the elements $e' \in E$
with $e' \leq e$. Then $M = \bigcup_{e \in E} M_e$.
Note that for each $e \in E$ the quotient
$$
M_e/\bigcup\nolimits_{e' < e} M_{e'}
$$
is either zero or generated by one element, hence has projective
dimension $\leq n$ by (3). By Lemma \ref{lemma-colimit-projective-dimension}
this means that $M$ has projective dimension $\leq n$.
\end{proof}

\begin{lemma}
\label{lemma-localize-finite-gl-dim}
Let $R$ be a ring. Let $M$ be an $R$-module.
Let $S \subset R$ be a multiplicative subset.
\begin{enumerate}
\item If $M$ has projective dimension $\leq n$, then $S^{-1}M$ has
projective dimension $\leq n$ over $S^{-1}R$.
\item If $R$ has finite global dimension $\leq n$, then
$S^{-1}R$ has finite global dimension $\leq n$.
\end{enumerate}
\end{lemma}

\begin{proof}
Let $0 \to P_n \to P_{n - 1} \to \ldots \to P_0 \to M \to 0$
be a projective resolution. As localization is exact, see
Proposition \ref{proposition-localization-exact},
and as each $S^{-1}P_i$ is a projective $S^{-1}R$-module, see
Lemma \ref{lemma-ascend-properties-modules},
we see that $0 \to S^{-1}P_n \to \ldots \to S^{-1}P_0 \to S^{-1}M \to 0$
is a projective resolution of $S^{-1}M$. This proves (1).
Let $M'$ be an $S^{-1}R$-module.
Note that $M' = S^{-1}M'$.
Hence we see that (2) follows from (1).
\end{proof}










